% =====================================================================
%  Chapter: System Architecture
%  Thesis: Integrating Open LLMs into the Jupyter Notebook
%          Reproducibility Pipeline -- Erisa Zaimi, TU Chemnitz
%
%  This chapter fixes component contracts, interfaces, and data
%  contracts at the level of detail settled during system design.
%  Where the later implementation (Chapter 5) confirmed, refined, or
%  corrected a design decision made here, that is noted explicitly in
%  place rather than silently rewritten, so the chapter also documents
%  how the design evolved.
% =====================================================================

\chapter{System Architecture}
\label{chap:architecture}

This chapter specifies the design of the LLM-based repair layer that
this thesis adds to the FAIR Jupyter reproducibility
pipeline~\cite{samuel2024gigascience}.
The existing pipeline ends at detecting and logging notebook execution
failures; even its containerised
evolution~\cite{samuel2025docker}
stops at the point of execution, recording the residual failure without
explaining or repairing it. The components described here attach at that
point and carry a dependency-related failure through explanation, repair,
re-execution, and structured logging.

Following the taxonomy of Yang et al.~\cite{yang2025aprsurvey},
the repair layer is realised as a \emph{procedural pipeline} rather than
an agentic framework: a fixed sequence of components driven by a
deterministic orchestrator. This choice is deliberate. Agentic loops
achieve strong results on complex, multi-file bugs but incur latency and
control overhead that scale poorly across a batch of dependency-error
notebooks; the interactive notebook-repair agents of Grotov et
al.~\cite{grotov2024untangling,grotov2024debug}
further assume a live kernel and a user in the loop, neither of which is
available in autonomous batch processing. The design goals follow from
the non-functional requirements: components are modular with an injected
configuration surface (NFR5), rely exclusively on open, locally
deployable models (NFR1), process the full target set without manual
intervention (NFR4), and are fully documented and version-controlled for
reproducibility (NFR3).

The scope of this chapter is the \emph{contracts} between components and
the \emph{data} they exchange. Prompt design and model selection are
described in \Cref{sec:m-explanation,sec:m-repair}; the internals of the
PyPI-grounded retrieval step are described in \Cref{sec:m-repair}. How each
component was actually built, including where implementation departed from
the contracts fixed here, is the subject of \Cref{chap:methodology} as a
whole.

% ---------------------------------------------------------------------
\section{Design Overview}
\label{sec:design-overview}

The system can be stated in one sentence: \emph{when a notebook has
failed because of a dependency error, read that failure, explain it in
plain language, derive a fix from live PyPI data, apply the fix and
re-run the notebook to check it, and record the result.} That sentence
is partitioned into five components, each realising one thesis
objective, sequenced by an orchestrator that realises pipeline
integration (O5):

\begin{table}[t]
  \centering
  \begin{tabular}{|l|l|p{7.2cm}|}
    \hline
    \textbf{Component} & \textbf{Objective} & \textbf{Responsibility} \\
    \hline
    ErrorClassifier  & FR1     & Confirm the failure is a fixable dependency error and identify the failing package (gatekeeper). \\
    \hline
    LLMExplainer     & O1      & Produce a plain-language explanation of the error. \\
    \hline
    RAGRepairAgent   & O2      & Derive a concrete fix, grounded in live PyPI metadata. \\
    \hline
    FixApplicator    & O3      & Apply the fix and re-execute the notebook to validate it. \\
    \hline
    ResultLogger     & O4, O6  & Persist the attempt to SQLite and enrich the knowledge graph. \\
    \hline
    Orchestrator     & O5      & Sequence the five components; expose the round and configuration switches. \\
    \hline
  \end{tabular}
  \caption{The five repair components and the orchestrator, each mapped
           to the objective it realises.}
  \label{tab:components-objectives}
\end{table}

% ---------------------------------------------------------------------
\section{Integration Point}
\label{sec:integration-point}

The containerised FAIR Jupyter pipeline~\cite{samuel2025docker}
executes the notebooks of each repository inside a per-repository Docker
container and writes one row per execution to the
\texttt{notebook\_executions} table, tagging each failure with an
\texttt{error\_category}. A pre-existing categorisation routine
(\texttt{summary.py}) already maps \texttt{ModuleNotFoundError} and
\texttt{ImportError} to \texttt{DEPENDENCY\_ERROR}. The repair layer
activates on exactly those rows: it performs a batch pass over the
\texttt{notebook\_executions} rows whose
\texttt{error\_category = DEPENDENCY\_ERROR}.

\paragraph{Re-execution substrate.}
A candidate fix is validated by re-running the notebook \emph{inside the
same Docker container the failure originated from}, top to bottom. Using
the originating environment ensures that a successful re-execution is
attributable to the applied fix rather than to an incidental difference
in the environment.

\paragraph{Context available at the hook-in point.}
Each target row exposes the following fields, which form the raw input to
the repair layer:
\texttt{error\_type}, \texttt{error\_message}, \texttt{error\_cell\_index},
\texttt{notebook\_id}, \texttt{repository\_id}, \texttt{notebook\_name},
\texttt{url}, \texttt{total\_code\_cells}, \texttt{executed\_cells}
(and \texttt{error\_category}).
Additionally, the cloned repository is available on disk, providing
\texttt{requirements.txt} and the \texttt{.ipynb} source; when the short
\texttt{error\_message} is insufficient, the full traceback is readable
from the legacy \texttt{executions.msg} column.

\paragraph{Working set.}
The committed Docker-run database contains \textbf{214}
\texttt{DEPENDENCY\_ERROR} rows: 172 \texttt{ModuleNotFoundError} and 42
\texttt{ImportError}. Dataset preparation (\Cref{sec:m-dataset}) resolved
the exact pip-fixable working set: \textbf{200 usable} rows and
\textbf{14 excluded} rows (10 system-library failures needing
\texttt{apt-get} rather than \texttt{pip}; 4 ambiguous local-module names
that cannot be safely mapped to a PyPI package). This is an exact count, not
an estimate; \Cref{chap:problem-dataset} gives the full breakdown, including
the further 13/187 development/evaluation split of the 200 usable rows. The
larger reproducibility figures from the complete conda
run~\cite{samuel2024gigascience} motivate the scale of the problem and are
reported in \Cref{chap:problem-dataset} as well, kept separate from this
Docker-pipeline working set. Should 200 targets prove insufficient for the
eventual full evaluation, the Grotov buggy-notebook
dataset~\cite{grotov2024untangling} remains a named fallback.

% ---------------------------------------------------------------------
\section{Component Architecture}
\label{sec:components}

\begin{figure}[t]
  \centering
  \resizebox{\textwidth}{!}{%
  \begin{tikzpicture}[
      font=\small,
      comp/.style={draw, rounded corners, fill=blue!6, minimum width=2.4cm,
                   minimum height=1cm, align=center, node distance=6mm},
      ext/.style={draw, rounded corners, fill=gray!10, minimum width=2.4cm,
                  minimum height=0.9cm, align=center, node distance=6mm},
      arr/.style={-{Latex[length=2.2mm]}, thick},
      darr/.style={-{Latex[length=2.2mm]}, thick, dashed},
    ]
    % five sequenced components
    \node[comp] (ec) {ErrorClassifier};
    \node[comp, right=of ec] (le) {LLMExplainer};
    \node[comp, right=of le] (ra) {RAGRepairAgent};
    \node[comp, right=of ra] (fa) {FixApplicator};
    \node[comp, right=of fa] (rl) {ResultLogger};

    % input above the chain
    \node[ext, above=14mm of ec] (input) {Docker pipeline\\\texttt{notebook\_executions}\\(dependency-related failure)};
    % PyPI above RAGRepairAgent
    \node[ext, above=14mm of ra] (pypi) {PyPI\\(live metadata)};
    % outputs well below ResultLogger, clear of the feedback curve
    \node[ext, below left=24mm and -12mm of rl] (sql) {SQLite\\\texttt{repair\_attempts}};
    \node[ext, below right=24mm and -12mm of rl] (kg) {FAIR Jupyter KG\\(RDF)};

    \draw[arr] (input) -- (ec);
    \draw[arr] (ec) -- (le);
    \draw[arr] (le) -- (ra);
    \draw[arr] (pypi) -- (ra);
    \draw[arr] (ra) -- (fa);
    \draw[arr] (fa) -- (rl);
    \draw[arr] (rl) -- (sql);
    \draw[arr] (rl) -- (kg);

    % bounded second-round feedback path (implemented, i7 - see
    % scripts/run_pipeline.py and sec:orchestrator)
    \draw[darr] (fa.south) .. controls +(0,-9mm) and +(0,-9mm) .. (ra.south)
      node[midway, below, font=\scriptsize] {bounded second round (implemented, i7; max.\ two rounds)};

    % orchestrator: dashed enclosure around the five sequenced components only
    \node[draw, dashed, rounded corners, inner sep=5mm, fit=(ec)(le)(ra)(fa)(rl)] (orch) {};
    % caption placed above LLMExplainer's column, clear of the input and PyPI arrows
    \node[font=\scriptsize, anchor=south] at (le.north |- orch.north)
      {Orchestrator (sequences the five components; implemented, i7)};
  \end{tikzpicture}%
  }
  \caption{Component architecture of the repair layer. The Docker pipeline's
           residual dependency failures enter at the left; the five
           components are sequenced by the orchestrator (dashed enclosure),
           implemented in \texttt{scripts/run\_pipeline.py} and validated by a
           small real pilot (\Cref{sec:orchestrator}); RAGRepairAgent queries
           PyPI for live metadata; ResultLogger writes both the SQLite
           benchmark table and the FAIR Jupyter knowledge graph. The dashed
           feedback arrow is the bounded second round, also implemented and
           capped at two rounds total. The full 187-row evaluation run over
           this pipeline remains future work (\Cref{chap:future-work}).}
  \label{fig:architecture-components}
\end{figure}

Each component is defined below by its responsibility and its input/output
contract. Implementation is deferred to \Cref{chap:methodology}.

\subsection{ErrorClassifier (FR1)}
\label{subsec:errorclassifier}
\textbf{Responsibility.} Decide whether a failure is an in-scope,
pip-fixable dependency error and, if so, extract the failing module and
its sub-type. It reuses the pipeline's existing \texttt{error\_category}
tag as a first filter and refines from \texttt{error\_type} and
\texttt{error\_message}.\\
\textbf{Input.} A \texttt{notebook\_executions} row (the fields listed in
Section~\ref{sec:integration-point}).\\
\textbf{Output.} A classification
$\{\texttt{failing\_module},\ \texttt{subtype},\ \texttt{scope\_status}\}$
where $\texttt{subtype} \in \{\texttt{missing\_package},\
\texttt{wrong\_version},\ \texttt{system\_library},\
\texttt{mapping\_unknown}\}$ and $\texttt{scope\_status} \in
\{\texttt{usable},\ \texttt{excluded}\}$.

\textbf{Explanation is not gated by scope.} \texttt{scope\_status} decides
whether RAGRepairAgent and FixApplicator run on a row, not whether
LLMExplainer does. Every classified row, including an \texttt{excluded} one,
is still passed to LLMExplainer: a researcher whose failure falls outside
the pip-only repair scope still deserves an explanation of what went wrong.
An earlier draft of this chapter placed one shared gate in front of both
LLMExplainer and RAGRepairAgent; \Cref{chap:problem-dataset} and
\Cref{chap:methodology} explain why that was corrected once the repair
component's implementation made the distinction concrete.

\subsection{LLMExplainer (O1)}
\label{subsec:llmexplainer}
\textbf{Responsibility.} Generate a plain-language explanation of the
error for a non-expert reader (NFR2).\\
\textbf{Input.} The error context (error type, message, failing module,
relevant cell).\\
\textbf{Output.} An explanation string and the identifier of the model
that produced it (\texttt{model\_used}), forwarded to the ResultLogger
for logging in \texttt{repair\_attempts.llm\_model}.\\
\emph{The prompt template and model are defined in
\Cref{chap:methodology}; this component's contract is fixed here, its
prompt internals are not.}

\subsection{RAGRepairAgent (O2)}
\label{subsec:repairagent}
\textbf{Responsibility.} Derive a concrete fix for the classified error.
Following PLLM~\cite{bartlett2025pllm},
the suggestion is grounded in live PyPI metadata so that version numbers
are retrieved rather than hallucinated.\\
\textbf{Input.} The classified record (error context, \texttt{failing\_module}),
the repository's requirements and import statements, the injected
model/prompt configuration, and -- in two-round mode -- the prior failed
attempt's new error as optional context.\\
\textbf{Output.} A \emph{fix object} (Section~\ref{subsec:fix-object}).\\
\emph{How PyPI is queried -- endpoints, fields consumed, import-name vs.
install-name resolution, and version selection -- is described in
\Cref{sec:m-repair}. This chapter fixes only the output contract (the fix
object); \Cref{sec:m-repair} also gives the exact field names the
implementation persists, which differ from the illustrative object shown in
Section~\ref{subsec:fix-object} below.}

\subsection{FixApplicator (O3)}
\label{subsec:fixapplicator}
\textbf{Responsibility.} Apply the fix object's command inside an
environment equivalent to the one the repository originally failed in, and
re-execute the notebook to validate it. The notebook is re-run \emph{top to
bottom}, not from the failing cell, because notebook statefulness means a
partial re-run cannot establish that the fix is correct~\cite{grotov2024debug}.\\
\textbf{Input.} A fix object, re-joined against the originating
notebook/repository record, and a rebuilt container equivalent to the one
the failure originally occurred in (\Cref{sec:m-fixapply} explains why this
is a rebuilt environment rather than a reused one, and why the join is
needed).\\
\textbf{Output.} An outcome
$\in \{\texttt{fixed},\ \texttt{still\_failing},\ \texttt{apply\_error}\}$,
together with the new error type/message if the notebook still fails.

\subsection{ResultLogger (O4, O6)}
\label{subsec:resultlogger}
\textbf{Responsibility.} Persist the complete attempt and enrich the
knowledge graph.\\
\textbf{Input.} The classification, explanation, fix object, outcome, and
run metadata.\\
\textbf{Output.} One \texttt{repair\_attempts} row
(Section~\ref{subsec:repair-table}) and, via the RML mapping, the
corresponding RDF triples (Section~\ref{subsec:kg}). The benchmark
dataset (O4) and the KG enrichment (O6) are both produced from this
single component.\\
\emph{This contract was built and validated by the result-logging and
knowledge-graph implementation (\Cref{sec:m-resultlogger}), including how
the four upstream JSONL outputs are actually joined -- a detail this
chapter leaves unspecified, since a join strategy is an implementation
concern, not a contract one -- and checked against real pilot data
(\Cref{sec:v-i6}).}

% ---------------------------------------------------------------------
\section{Orchestration and Control Flow}
\label{sec:orchestrator}

The orchestrator iterates over the target rows and runs the components in
order: classify, explain, repair, apply, log. Three properties are fixed
at this level.

% TODO-FULL-RUN: update after full orchestrated dataset execution
\fullrunupdate{The orchestrator described below is now implemented
(\texttt{scripts/run\_pipeline.py}) and validated by a small real pilot run
against real Ollama, PyPI, and Docker, recorded in the accompanying
implementation documentation (\texttt{docs/pipeline-orchestrator.md}) --
not the full 200-row usable set. Each of ErrorClassifier through ResultLogger exists,
is independently tested, and is called by the orchestrator exactly as
described below, including the bounded two-round feedback mode described
later in this section. What remains is the full 200-row usable-set
run itself, the resulting one-round-vs-two-round comparison, and full
benchmark/knowledge-graph population from it -- \Cref{chap:future-work}.
Result logging and knowledge-graph enrichment were already implemented
before the orchestrator (\Cref{sec:m-resultlogger}) and are unchanged by
it; the orchestrator only calls them.}

\paragraph{Round switch.}
The default is a single repair round. In two-round mode, if the outcome is
\texttt{still\_failing}, the new error is fed back to the RAGRepairAgent
once more and the resulting fix is re-applied before the attempt is
logged. This feedback structure follows the re-execution-driven repair of
ChatRepair~\cite{xia2024chatrepair};
the one- versus two-round comparison is the subject of the ablation study.

\paragraph{Failure isolation (NFR4).}
Every attempt is logged, even when a component itself errors -- malformed
LLM output, an unreachable PyPI, or a re-execution timeout. Such cases are
caught and recorded with the appropriate \texttt{outcome}; a single
problematic notebook never aborts the batch.

\paragraph{Configuration surface (NFR5).}
The model, the prompt strategy, and the round count are injected per run
rather than hard-coded, so they can be swapped without changing component
code. The same injection points drive the model and prompt comparisons in
the evaluation.

\begin{verbatim}
for row in notebook_executions where error_category = 'DEPENDENCY_ERROR':
    cls = ErrorClassifier(row)
    if not cls.is_dependency_error:
        continue
    explanation = LLMExplainer(row, cls)            # O1
    fix         = RAGRepairAgent(row, cls)          # O2  (PyPI-grounded)
    outcome     = FixApplicator(fix, container)     # O3  (re-run top-to-bottom)
    log(row, cls, explanation, fix, outcome, round=1)

    if two_round and outcome == 'still_failing':    # second round
        fix2     = RAGRepairAgent(row, cls, prev_error=outcome.new_error)
        outcome2 = FixApplicator(fix2, container)
        log(row, cls, explanation, fix2, outcome2, round=2)
\end{verbatim}
% LISTING: switch the block above to lstlisting/algorithm if available.

% ---------------------------------------------------------------------
\section{Data Contracts}
\label{sec:data-contracts}

\subsection{Fix Object}
\label{subsec:fix-object}

The RAGRepairAgent emits the \emph{proposed} fix as structured JSON,
before application. Version~1 is pip-only: \texttt{action} is one of
\texttt{install}, \texttt{pin\_version}, or \texttt{none}; \texttt{version}
is \texttt{null} for a plain install; and \texttt{command} is exactly what
the FixApplicator runs inside the container. Notebook-code edits
(\texttt{replace\_import}) are deferred to future work -- in practice many
renamed-import cases are still resolved by pinning an older version where
the import was valid, which remains a pip fix.

\emph{This is the object's conceptual shape, fixed at design time.} The
implemented object (\Cref{sec:m-repair}) keeps the same four meaningful
values -- action, resolved install name, version, and the evidence behind
them -- but names and nests the fields differently
(\texttt{final\_action}, \texttt{final\_install\_name}, and so on, with
retrieval evidence in a separate \texttt{retrieval\_result} block) once the
validation logic in \Cref{sec:m-repair} was actually built. The example
below should be read as the contract's intent, not as the literal persisted
record.

\begin{verbatim}
{
  "action": "pin_version",
  "import_name": "scanpy",
  "install_name": "scanpy",
  "version": "1.9.3",
  "command": "pip install scanpy==1.9.3",
  "rationale": "Notebook needs an API present only up to 1.9.x; pinning the
                last compatible release.",
  "pypi_evidence": {
    "latest_version": "1.10.1",
    "chosen_version": "1.9.3",
    "requires_python": ">=3.9"
  }
}
\end{verbatim}
% LISTING: JSON

\subsection{The \texttt{repair\_attempts} Table}
\label{subsec:repair-table}

The repair layer writes to a single new SQLite table,
\texttt{repair\_attempts}, which serves a dual purpose: it is the live
validation log (O3) and the reusable benchmark dataset (O4). It records
\emph{one row per attempt}: a two-round repair produces two rows, and
re-running the same targets under a different model or prompt produces
further rows. The error data is never copied -- each row joins back to the
originating failure via \texttt{notebook\_execution\_id}. Filtering on
\texttt{round}, \texttt{llm\_model}, \texttt{prompt\_strategy}, and
\texttt{run\_id} is what powers the ablation and the model/prompt
comparisons.

\begin{verbatim}
CREATE TABLE repair_attempts (
  id                      INTEGER PRIMARY KEY,
  notebook_execution_id   INTEGER NOT NULL,   -- FK -> notebook_executions.id

  -- from ErrorClassifier
  failing_module          TEXT,               -- e.g. "scanpy"
  subtype                 TEXT,               -- missing_package | wrong_version

  -- from LLMExplainer (O1)
  explanation             TEXT,

  -- from RAGRepairAgent -- the proposed fix (O2)
  action                  TEXT,               -- install | pin_version | none
  install_name            TEXT,               -- resolved PyPI name
  version                 TEXT,               -- version to pin, or NULL
  command                 TEXT,               -- exact command run
  rationale               TEXT,
  pypi_evidence           TEXT,               -- JSON: versions seen, requires_python

  -- from FixApplicator -- the outcome (O3)
  outcome                 TEXT,               -- fixed | still_failing | apply_error
  new_error_type          TEXT,
  new_error_message       TEXT,

  -- run metadata (for experiments)
  llm_model               TEXT,
  prompt_strategy         TEXT,               -- for the prompt-strategy comparison
  round                   INTEGER,            -- 1 or 2 (the ablation)
  run_id                  TEXT,               -- groups one run / config
  created_at              TEXT,

  FOREIGN KEY (notebook_execution_id) REFERENCES notebook_executions(id)
);
\end{verbatim}
% LISTING: SQL

The result-logging implementation built this table exactly as specified here:
\texttt{scripts/result\_logger.py}'s own \texttt{CREATE TABLE} statement
reproduces this nineteen-column schema unchanged, the one data contract in
this chapter that implementation did not have to revise once it was actually
built. \Cref{sec:m-resultlogger} describes how each of the eighteen
non-\texttt{id} columns is populated from the four upstream JSONL files,
including how a \texttt{notebook\_execution\_id} that legitimately repeats
across separate attempts is kept as separate rows rather than collapsed.

\subsection{Knowledge-Graph Enrichment}
\label{subsec:kg}

The FAIR Jupyter knowledge graph~\cite{samuel2024fairkg}
is not written by hand; it is generated from per-table RML mappings
(\texttt{*.rml.ttl}) run by the existing build script
(\texttt{run\_fairjupyter\_kg.sh}), using Morph-KGC as the mapping engine.
Objective~O6 adds a single mapping, \texttt{repair\_attempts.rml.ttl},
modelled on the existing \texttt{mapping/rml\_mapping/executions.rml.ttl}. It
re-expresses each \texttt{repair\_attempts} row as triples and attaches them
to the notebook node already present in the graph, so that repair history
becomes queryable alongside the existing reproducibility data. This mapping
follows the existing FAIR Jupyter knowledge-graph architecture -- CSV-based
RML mappings run through Morph-KGC -- rather than introducing a separate
RDF-generation pipeline of its own.

The vocabulary strategy follows the earlier design notes: \emph{reuse} existing
\texttt{repr:} terms where they already fit (e.g.\ \texttt{repr:exception}),
\emph{add} new \texttt{repr:} terms for repair-specific facts, and use
PROV-O~\cite{provenance2013}
for the activity/agent/time layer, representing each repair as a
\texttt{prov:Activity}. This serialisation of repair provenance directly
serves the Reusable principle of FAIR~\cite{wilkinson2016fair} and the
execution-environment provenance emphasised by
FAIR4RS~\cite{barker2022fair4rs}.
The base IRI is \texttt{repr:}~$=$~\url{https://w3id.org/reproduceme/} and
the serialisation is Turtle (\texttt{.ttl}).

\begin{verbatim}
repr:notebook_512  repr:hadRepairAttempt  repr:repairattempt_001 .

repr:repairattempt_001
    a                       prov:Activity, repr:RepairAttempt ;
    repr:exception          "ModuleNotFoundError" ;   # reuse existing term
    repr:failingModule      "scanpy" ;                # new term
    repr:appliedFix         "pip install scanpy==1.9.3" ;
    repr:fixOutcome         "fixed" ;
    prov:wasAssociatedWith  repr:codellama_13b ;       # which model did it
    prov:endedAtTime        "2026-07-15T14:22:00Z"^^xsd:dateTime .
\end{verbatim}
% LISTING: Turtle. NOTE: this is the mapping's design-time intent, not its
% literal implemented shape. The implemented mapping (\Cref{sec:m-resultlogger})
% keeps the same reuse-vs-add vocabulary strategy and the same
% prov:endedAtTime typing, but does not use repr:appliedFix or
% prov:wasAssociatedWith, and reuses repr:exception/repr:msg for the
% post-repair error rather than the original one, which the existing
% executions mapping already records; the exact predicate names actually
% generated are listed there.

% ---------------------------------------------------------------------
\section{Scope and Forward References}
\label{sec:arch-scope}

This chapter fixes the component contracts, the orchestration logic, and
the three data contracts (fix object, \texttt{repair\_attempts} table, KG
mapping). The following are covered elsewhere in this thesis, or remain
future work:

\begin{itemize}
  \item Prompt templates and open-model selection -- \Cref{sec:m-explanation,sec:m-repair}.
  \item PyPI retrieval internals: endpoints, consumed fields, import-name
        vs.\ install-name resolution, and version selection --
        \Cref{sec:m-repair}.
  \item Component implementation, the real data shapes produced, and the
        design deviations discovered while building them, including how
        ResultLogger joins the four upstream outputs and how the
        knowledge-graph mapping was implemented -- \Cref{chap:methodology};
        their validation against real pilot data -- \Cref{chap:validation}.
  \item Full end-to-end orchestration and the bounded two-round feedback
        mode -- \textbf{implemented} (\texttt{scripts/run\_pipeline.py});
        the full 200-row usable-set run and the resulting
        one-round-vs-two-round comparison remain future work -- see
        \Cref{chap:future-work}.
\end{itemize}
